
\section*{Ventana N57: cierre Bockstein--Hensel}
Para una columna:
\[
(p,e,f)\in\{0,1,2\}^3.
\]
\[
q=(p+e+f)\bmod3,
\qquad
c=\frac{p+e+f-q}{3}.
\]
Canal firmado:
\[
a=(p+e-f)\bmod3,
\qquad
h=\frac{p+e-f-a}{3}.
\]
Con:
\[
\kappa(x,y)=\frac{x+y-((x+y)\bmod3)}{3},
\quad
\beta(x,y)=\frac{x-y-((x-y)\bmod3)}{3},
\]
\[
\boxed{c=\kappa(p,e)+\kappa((p+e)\bmod3,f)}
\]
\[
\boxed{h=\kappa(p,e)+\beta((p+e)\bmod3,f)}.
\]
Además:
\[
\delta\kappa=0.
\]
Los residuos \(q,a\) son afines, pero:
\[
\boxed{c,h\text{ no son afines: son la obstrucción de carry.}}
\]
Ecuaciones globales:
\[
\boxed{X_{t+1}^+=\frac{X_t^+A-q_t}{3}-c_t}
\]
\[
\boxed{X_{t+1}^\sigma=\frac{X_t^\sigma A-a_t}{3}-h_t}.
\]
Terminal:
\[
q=122120,\qquad c=000110,\qquad r=221.
\]
Dictamen:
\[
\boxed{P_m\text{ no es primitivo; el carry Bockstein--Hensel es la ley.}}
\]
\[
\boxed{\text{blanco restante: derivar }x_t^\chi,u_t^\chi\text{ desde APP/TPK.}}
\]
